% ============================================================
% Introduction — draft v3
% All citation keys resolve against references.bib (v. 2026-10-16d).
% Requires: natbib, enumitem.
%
% OPEN ITEMS FOR THIS SECTION (each is also marked TODO at its place):
%  [I1] Appendix app:traffic (production tile sizes) does not exist yet. Write it from
%       real-world-map-traffic.md: method, limits, and the tile-size table only
%       (no page shares, no Aviasales MLT estimate).
%  [I2] Yandex median tile size (3.8 KB) may mix raster PNG and vector tiles; recompute
%       the median for vector tiles only and update the range "0.1 and 3.8 KB".
%  Extended experiment: content-normalized results and address totals are generated
%  from the five-variant Rust cohort; historical three-variant counts do not apply.
%  [I5] Section labels (sec:results-granularity, sec:results-density,
%       sec:results-streams, sec:results-bpv, sec:results, sec:discussion,
%       sec:threats) must exist in the later sections.
%  [I6] DONE: novelty claim narrowed after literature/technical review; prior work observes
%       larger gains for naturally large tiles but does not run a controlled spatial-granularity sweep.
%  [I7] Re-check OSM node count and Overture figure against their sources on the
%       submission date (both cited with access dates).
%  [I8] 2026-10-17: motivation added (display format, motivating settings, three
%       observations). New keys space_apps, nasa_mars_trek_wmts in references.bib.
%       Update the abstract's first sentence to match (see suggestion in chat).
%       The motivation is closed in Discussion ("Back to the motivating settings")
%       and in the Conclusion; keep the three in sync if the wording changes.
%  [I9] Consider an Acknowledgments/competing-interests line on the commercial
%       fleet-management system (inspiration only; no data or code shared).
% ============================================================

\section{Introduction}
\label{sec:introduction}

Interactive maps are drawn on the client from tiles. The map is cut into a pyramid of
zoom levels and, at every level, into a regular grid; each tile is encoded, stored and
transferred independently, and the client requests the tiles that cover its current view
and styles and renders their content itself \citep{osm_slippy_tiles, ogc2022tms}.
Vector tiles are therefore a display format: they carry geometry and attributes already
selected and generalized for drawing at one scale, not the source data from which they
were derived. The compactness of this representation matters well beyond download time.
It determines the storage cost of map archives that cover the whole planet,
the feasibility of offline maps on mobile devices and in low-connectivity settings,
and the volume of data that clients must transfer and decode.
These pressures grow with the underlying data. OpenStreetMap passed ten billion
nodes in August 2025 \citep{osm_wiki_stats}, and as early as 2024 the Overture Maps
release combined 2.35 billion building footprints from OpenStreetMap, Esri Community
Maps and machine-learning extractions by Microsoft and Google \citep{overture2024alpha}.

This study grew out of practical work with maps of this kind. In a commercial
fleet-management system for ground robots that operate in remote areas with little or no
connectivity, the on-site server, a computer of the Raspberry Pi class, also serves maps
to the operator. The operator views the area with custom layers that depend on the
robot's task, such as agricultural work, clearing vegetation under power lines or snow
removal. A similar need arises in a student project for the NASA International Space
Apps Challenge \citep{space_apps} on the visualization of planetary surfaces, which NASA
itself distributes as web map tiles \citep{nasa_mars_trek_wmts}. Both settings display
custom thematic layers, which are often sparse, from stored tile archives, consider MLT as
the encoding, and are constrained by storage, link capacity or computing power. Neither
system uses the results reported here, and no data from them enter the experiment. They
raised the question that this paper studies in isolation: how far the tiling scheme,
rather than the encoding format, limits how compactly such maps can be stored and
delivered.

The dominant encoding for vector tiles has been the Mapbox Vector Tile (MVT) format
\citep{mapbox_mvt_spec_2_1}, which stores each feature as a Protocol Buffers message
that holds its geometry and its attributes; in this sense MVT is feature-oriented.
MapLibre Tiles (MLT) is a recently proposed successor that reorganizes tile content
into columns: geometry and each attribute are stored in separate streams and compressed
with lightweight encodings such as dictionary, run-length and delta encoding and
integer bit-packing \citep{tremmel2025mlt, maplibre_mlt_v1}.
Column-oriented layouts with such encodings are well established in analytical
databases and data lakes \citep{abadi2006compression, zeng2023columnar, kuschewski2023btrblocks},
and spatial variants have been proposed for geospatial data lakes \citep{saeedan2022spatialparquet}.

The efficiency of columnar encoding, however, depends on how much data each
independently encoded unit contains. Every stream carries structural information, and
lightweight encodings have more opportunity to exploit regularity when a stream contains
more values. In columnar storage this dependence is well known and is controlled
explicitly: Apache Parquet lets the writer choose the size of row groups and data pages,
and its documentation notes that larger pages carry less header overhead and recommends
large row groups \citep{parquet_configurations}.

We use \emph{spatial tile granularity} in a specific sense: the geographic area and
corresponding content grouped into one independently encoded tile at an approximately
fixed logical level of detail. This is distinct from \emph{coordinate extent}, the
integer coordinate grid used inside a tile; from \emph{encoded tile size}, measured in
bytes; and from the rendered tile size in screen pixels. In a conventional pyramid the
spatial granularity is imposed by the tile matrix rather than chosen by the encoding
format. The MVT specification deliberately keeps bounds and projection outside the tile
\citep{mapbox_mvt_spec_2_1}; the quadtree pyramid, in which each zoom level doubles the
linear resolution, comes from tiling conventions shared by servers and renderers
\citep{osm_slippy_tiles, ogc2022tms}. Generalization policies keep tiles small as well,
since production tilers drop or simplify features to bound tile size.
In a measurement of seven production web maps (Appendix~\ref{app:traffic}),
the vector tiles we observed were typically only a few kilobytes on the wire,
with per-product medians between 0.1 and 3.8~KB.
% TODO [I1][I2]: write the appendix; vector-only median for Yandex.

Seen from the display client, this fixed unit appears smaller than necessary. Three
observations from the settings above led us to the present question. First, the rendered
tile size is fixed: MapLibre GL JS draws vector tiles at 512 CSS pixels
\citep{maplibre_gljs_vectorsource}, so at an integer zoom level a viewport of
$2560\times1600$ CSS pixels, the one used in Appendix~\ref{app:traffic}, intersects 20 to
30 tiles at once. The client thus already holds a block of several tiles for a single
view. Second, typical map tasks, such as surveying an area to choose a site for
equipment, span several screens and several zoom levels, so a session requests many more
tiles than one view: in Appendix~\ref{app:traffic}, a session of two pans and one zoom
step loaded between about 270 and 770 tiles in six of the seven products. These two
observations suggest that a coarser encoding unit need not cost much additional transfer,
although we do not measure that cost here (Section~\ref{sec:discussion}). Third, and most
important, the density of map content varies widely between tiles, whereas at a given
zoom level every tile covers the same area. Dense urban tiles are already large, while
tiles of sparse thematic layers hold a handful of features and consist largely of
per-tile structure. Because the unit is fixed by the tile matrix, an encoder cannot
enlarge it where the content is sparse. In our own map applications we have worked around
this limit with overzoom and by combining several tilesets with different maximum zoom
levels; such workarounds reduce the number of tiles but do not change the unit of
compression in a controlled way (Section~\ref{sec:bg-overzoom}).

That larger blocks of data compress better is unsurprising in itself.
Less obvious is how strongly a columnar tile format depends on spatial encoding-unit
granularity compared with a feature-oriented one, and whether the effect survives
general-purpose transport compression such as gzip \citep{rfc1951, rfc1952}. The
original MLT evaluation reports compression up to three times better than MVT across
tilesets and more than six times better for some large tiles
\citep{tremmel2025mlt}. MapLibre's current planet-scale benchmarks make the association
more explicit: naturally large tiles obtain larger MLT gains than median tiles
\citep{maplibre_mlt_benchmarks}. Those results stratify existing tiles by observed
encoded size; they do not manipulate the spatial boundary of the independently encoded
unit while holding logical level of detail approximately fixed.

Prior work has also varied vector-tile partitioning, zoom, pyramid depth and tile
content, but these interventions either change level of detail together with spatial
scale, optimize content within a fixed tile, or change pyramid organization rather than
the compression-unit boundary \citep{li2018geotile, dufilie2014feathered,
bahadori2026hifive, tippecanoe, planetiler_1574}. To our knowledge, no previous study has
systematically varied the spatial granularity of independently encoded vector tiles
under approximately matched-content, matched-LoD conditions and quantified the resulting
raw and per-tile-gzip size of column-oriented MLT relative to feature-oriented MVT.


In this paper we distinguish fixed-content aggregation from regeneration of larger
tiles. Aggregated-4 and Aggregated-16 regroup the exact Standard feature multiset
into larger encoding units; Dense-4 and Dense-16 regenerate from the same normalized
sources at matched logical zoom. Both interventions are evaluated as MVT, MLT v1 and
experimental MLT v2 using one pinned Rust implementation.

We ask whether fixed-content aggregation reduces bytes, how regeneration changes
that effect, whether it persists across MLT versions, and which stream-level and
content-normalized measurements are consistent with the outcome. Our contributions
are an exact typed-multiset aggregation procedure; a geographic block-matched,
additive decomposition of byte changes; and a reproducible five-variant,
three-format comparison on five heterogeneous datasets. Source-feature provenance
supports counts of unique normalized records and repeated instances. The original
Python-produced corpus is retained as historical reference.

Section~\ref{sec:background} describes the formats and related work;
Section~\ref{sec:method} defines the controls and measurements;
Section~\ref{sec:results} presents the expanded results;
Sections~\ref{sec:discussion} and~\ref{sec:threats} discuss their interpretation.
